% Figure 12 -- graphic form of the Seiberg-Witten curve
\begin{tikzpicture}[x=1pt,y=1pt,scale=1.08,transform shape,>=stealth,line width=0.5pt]
  % ---- axes ----
  \draw (52.4,8.6) -- (52.4,175.3);
  \draw (4.6,56.4) -- (230.5,56.4);
  % ---- Newton-polygon-like outline ----
  \draw (99.6,56.0) -- (99.6,139.7) -- (147.5,115.5) -- (147.2,56.0);
  \draw (147.5,115.5) -- (206.7,56.2);
  % ---- k1, k2 brackets ----
  \draw (99.4,50.3) -- (117.9,50.4) -- (123.7,44.1) -- (129.7,50.6) -- (141.6,50.3);
  \draw (146.8,50.3) -- (171.3,50.4) -- (177.1,44.1) -- (183.0,50.6) -- (206.8,50.3);
  % ---- h1, h2 markers ----
  \draw (93.9,91.6) -- (81.9,139.8);
  \draw (93.9,92.3) -- (82.0,62.0);
  \draw (141.3,116.0) -- (135.3,85.7);
  \draw (141.3,62.0) -- (135.2,86.4);
  % ---- distinguished points ----
  \fill (82.0,145.3) circle[radius=1.85pt];
  \fill (70.1,151.2) circle[radius=1.85pt];
  % ---- labels ----
  \node[inner sep=0pt] at (181.25,35.6) {\fontsize{9}{11}\selectfont$k_1$};
  \node[inner sep=0pt] at (121.95,35.6) {\fontsize{9}{11}\selectfont$k_2$};
  \node[inner sep=0pt] at (127.95,88.9) {\fontsize{9}{11}\selectfont$h_1$};
  \node[inner sep=0pt] at (74.55,94.85) {\fontsize{9}{11}\selectfont$h_2$};
\end{tikzpicture}
